\section{Conclusions}\label{sec:conclusion}
The empirical contribution is an audit of source completeness linked to matched-origin price and interval evaluation, supported by strong simple benchmarks. Historical results suggest that missing reports can change the apparent calibration of a forecasting pipeline. They motivate checking data provenance before attributing error changes to market drift; they do not supply a new forecasting algorithm, new conformal theory or causal identification of reporting errors.

The clean-source release rebuilds both cleaning variants and reruns the smaller, separately calibrated study. Its observed average improvements against persistence are reported separately from the historical full-size model comparisons. Approximate long-block and HAC sensitivities assess, but do not resolve, dependence assumptions. Forecast rows are not independent experimental replicates; limited seasons, shared shocks, changing exposure and training uncertainty remain material limitations. A good historical pooled MAE does not imply improvement in every market or profitable farm decisions.

LightGBM's prototype adoption remains exploratory. Neither the release study nor the earlier sensitivity independently confirms its advantage over newly refitted Spark alternatives. A human statistician's review, immutable large-input archival, independent execution and genuinely future-issued evaluation remain outstanding. These are prerequisites for a stronger submission recommendation, not reasons to add an unmotivated neural model.

\paragraph{Data and code availability.}
The repository contains source, configurations, tests, experiment manifests and compact result files. The release fitting-source tag is \texttt{agdata-}\allowbreak\texttt{research-}\allowbreak\texttt{source-v1}; retained build records bind pre/post source, inputs and executed JAR hashes to commands and feature outputs. A local source archive and environment inventory accompany this bounded reproduction; no anonymised journal package or public immutable large-data archive is claimed supplied. Large data, model and prediction files remain outside Git and require the retained bytes matching the manifests. Later API responses may differ after source revisions.
The original full-size experiment recorded a dirty tree and a parent commit predating the forecasting code, so its exact execution state remains unproven. That historical gap is not repaired by the release and those full-size principal tables have not been rerun end-to-end here. The fixed historical panel is a hashed external input rather than a reproduced selection procedure.

Sources are the public Kaggle archive and official data.gov.in resource \texttt{35985678-}\allowbreak\texttt{0d79-}\allowbreak\texttt{46b4-}\allowbreak\texttt{9ed6-}\allowbreak\texttt{6f13308a1d24} (a key is required). Table-to-artifact mappings, commands and evidence classes are documented in \path{docs/VERIFIED_METRICS.md}, \path{docs/RELEASE_PROTOCOL.md} and \path{docs/EVIDENCE_LEDGER.md}. IJF is the intended empirical forecast-evaluation venue; this working draft is not submitted, and journal-template compliance, author details and declarations require verification.
